\section{Selecting split lines}\label{sec:selection}

The candidate graph supplies many choices in each slot; we must now
choose one vertex from every slot without creating a conflict.  We need
two further properties of this choice.  First, an event prescribing the
chosen vertices in several slots must have small probability.  Such an
event is called a \emph{cylinder}, and its bound will allow us to count
exceptional pencil configurations.  Second, the chosen split lines must
leave a positive proportion of the initially available points on every
affine line.

We prove these properties together in a graph statement whose bounds are
uniform over all inputs satisfying the hypotheses.  The cylinder bound
includes the event that the procedure succeeds; we do not condition the
individual choices on eventual success.

\begin{proposition}[Selection with bounded cylinder probabilities]
\label{sel:selection}
Fix constants $c_*,C_*>0$ and $0<c_0\le1$, and put $\beta=1/8$.
For each sufficiently large integer $q$, let a finite graph have its
vertex set partitioned into $n=2(q+1)$ slots $Q_i$, with no edge within
a slot.  Write $N(v)$ for the neighborhood of a vertex.  Suppose that
\begin{align}
 c_*q\le |Q_i|\le q,&\qquad |N(v)|\le C_*q,
 \qquad |N(v)\cap Q_i|\le q^\beta,\label{sel:graph-bounds}\\
 |J(v)|\le q^\beta,&\qquad
 |N(v)\cap N(w)|\le q^\beta\quad(w\notin J(v)),\label{sel:exceptional}
\end{align}
where every $w\in J(v)$ is either adjacent to $v$ or in its slot.
There are constants $K_*,\rho>0$, depending only on $c_*,C_*,c_0$,
and a random procedure which either fails or returns an independent
transversal $(V_i)_{i\in I_0}$, where $I_0$ is the set of slots and
$V_i\in Q_i$.  If $\mathcal A$ denotes its event of returning a
transversal, then
\begin{equation}\label{sel:success}
 \Prb(\mathcal A)=1-o(1).
\end{equation}
For any $s$ distinct slots $i_1,\ldots,i_s$ and any prescribed vertices
$v_j\in Q_{i_j}$,
\begin{equation}\label{sel:cylinder}
 \Prb\bigl(\mathcal A, V_{i_j}=v_j\ (1\le j\le s)\bigr)
 \le (K_*/q)^s.
\end{equation}

Suppose also that the slots are indexed by $(g,k)$, with two slots for
each direction $g$ of an affine plane of order $q$.  Associate to each
vertex in slot $(g,k)$ a line of direction $g$, injectively within
that slot.  For every affine line $\ell$, let $T_\ell\subseteq\ell$
be a fixed set of at least $c_0q$ points.  From $T_\ell$ remove every
point on a selected line different from $\ell$, and call the remaining
set $T_\ell^{\rm fin}$.  Then the same procedure satisfies
\begin{equation}\label{sel:occupancy-conclusion}
 \Prb\bigl(\mathcal A, \min_\ell |T_\ell^{\rm fin}|<\rho q\bigr)
 =o(1).
\end{equation}
The errors in \eqref{sel:success} and
\eqref{sel:occupancy-conclusion} are uniform in all the stated input
data.  By the injectivity within a slot, \eqref{sel:cylinder} also
holds for prescribed candidate lines.
\end{proposition}

The selection has two stages.  In the main stage we repeatedly choose a
slot and a vertex, delete conflicts, and thin the remaining choices so
that every vertex has the same survival probability.  Domain sizes then
decrease at a common rate.  Remaining degrees decrease faster, because
the chosen slot also disappears.  We show that this difference leaves a
sparse enough graph to complete the last slots by the local lemma.
Partial random selection followed by local-lemma completion also appears
in Loh and Sudakov~\cite[Section~4]{LohSudakov2007}. Here we additionally
track occupancy on each affine line and prove a cylinder bound for the
complete returned transversal.

The main technical issue is concentration of the remaining degrees.
A choice in $J(v)$ can delete many neighbors of $v$, but it also kills
$v$ or removes its slot.  We will use this implication to stop tracking
$v$ after such a choice, without losing any degree estimate needed for
a subsequent step.

\subsection{The main selection process}

We first specify the main stage, including the events on which it
aborts.  Choose the constants $P_*,s_{\min},B_*,\rho,a$ in the
displayed order, and define the $q$-dependent parameters
$\alpha,h_0,t_{\max}$ alongside them:
\begin{align}
 P_*&=4C_*/c_*+4,& \alpha&=P_*/q,&
 s_{\min}&=\exp(-6P_*),\label{sel:constants-one}\\
 B_*&=2/(s_{\min}c_*),&
 \rho&=c_0\exp(-6B_*)/4,\label{sel:constants-two}\\
 0<a&<\min\{1,\rho/2,c_*/(100(C_*+1))\},&
 h_0&=\lceil aq\rceil,& t_{\max}&=n-h_0.\label{sel:constants-three}
\end{align}
All subsequent assertions concern sufficiently large $q$.  The main
stage runs for at most $t_{\max}$ steps, leaving $h_0$ slots if it
reaches completion.
At time $t$, before the next step, let $I(t)$ be the remaining slots,
$h(t)=n-t$, and $Q_i(t)$ their domains of vertices still alive.  Put
\begin{equation}\label{sel:scales}
 L_i(t)=|Q_i(t)|,\qquad
 s(t)=(1-\alpha)^t,\qquad b(t)=s(t)h(t)/n.
\end{equation}
Initially all slots and vertices are present.  The factors $s(t)$ and
$b(t)$ will be the respective scales for domain sizes and remaining
degrees.  Throughout $0\le t\le t_{\max}$,
\begin{equation}\label{sel:scale-lower}
 s_{\min}\le s(t)\le1,\qquad s_{\min}a/3\le b(t)\le1.
\end{equation}
Indeed $n\le3q$, $h(t)\ge aq$, and
$\log(1-P_*/q)\ge-2P_*/q$ for large $q$.

Each main step begins with two checks.  First, abort if some remaining
domain satisfies
\begin{equation}\label{sel:size-barrier}
 L_i(t)<s(t)c_*q/2.
\end{equation}
For every alive vertex $v$ in a remaining slot define
\begin{equation}\label{sel:degree-hazard}
 D_v(t)=\sum_{i\in I(t)}|N(v)\cap Q_i(t)|,\qquad
 r_v(t)=\frac1{h(t)}\sum_{i\in I(t)}
            \frac{|N(v)\cap Q_i(t)|}{L_i(t)}.
\end{equation}
Here $r_v(t)$ is the probability that a uniform choice from a uniformly
chosen remaining slot is adjacent to $v$.  The second check aborts if
any such $r_v(t)>\alpha$.

If both checks pass, choose a uniform slot $G\in I(t)$ and then a
uniform vertex $W\in Q_G(t)$.  Record $W$ as the choice for $G$ and
delete all its neighbors.  For every vertex $u$ alive before this step
and not adjacent to $W$, independently delete it with probability
\begin{equation}\label{sel:thinning-rate}
 \lambda_u=\frac{\alpha-r_u(t)}{1-r_u(t)}.
\end{equation}
Compute all rates before choosing $G,W$.  The pre-step checks ensure
$0\le\lambda_u\le\alpha<1$.  For the estimates below, perform this
thinning hypothetically in slot $G$ as well: the recorded choice $W$
is retained regardless of its hypothetical deletion.  Finally remove
slot $G$.  Thus, before slot removal, every previously alive vertex
has the same conditional survival probability:
\begin{equation}\label{sel:equal-survival}
 (1-r_u(t))(1-\lambda_u)=1-\alpha.
\end{equation}

A final check takes place after the step.  Count the thinning deletions
in each domain present at its start and in each set
$N(v)\cap\bigcup_{i\in I(t)}Q_i(t)$ then counted by a $D_v(t)$.
Abort if any such count exceeds $q^\beta$, including hypothetical
thinning in $G$.  This check limits the size of the increments used
in the concentration argument.

The probability of a post-step abort is small uniformly over all
histories that pass the pre-step checks.  Conditional on such a history
and on $G,W$, each tested count is a sum of independent Bernoulli variables
with total mean at most
$\mu_*=P_*\max\{1,C_*\}$.
For such a sum $Z$ and any integer $k\ge1$,
\begin{equation}\label{sel:bernoulli-tail}
 \Prb(Z\ge k)\le \E\binom Zk\le\mu_*^k/k!
 \le(e\mu_*/k)^k.
\end{equation}
The middle inequality follows by expanding the factorial moment and
bounding the sum of products of $k$ distinct success probabilities
by $\mu_*^k/k!$.  With $k=\lfloor q^\beta\rfloor+1$, the final
bound, denoted $\varepsilon_q$, is smaller than $q^{-m}$ for every
fixed $m$ once $q$ is sufficiently large.  There are $O(q^2)$ tested
sets per step and at most $3q$ steps.  Thus any post-step abort has
probability $O(q^3\varepsilon_q)=o(1)$.

\subsection{Concentration with stopping and censoring}

We now show that the main stage reaches the last $h_0$ slots with
large domains and small remaining degrees.  Let $\mathcal F_t$ be the
original history just before the next random slot and vertex are
chosen.  All conditional expectations in this subsection use this
filtration; in particular, the next slot is still uniform among the
remaining slots.

We use auxiliary processes that agree with the quantities of interest
until they are stopped.  Stopping means holding the last value fixed.
A pre-step abort produces no increment.  A performed step contributes
its specified increment before any post-step stop.  When a large
decrement is censored below, only the auxiliary update is replaced;
the actual selection procedure is unchanged.

The concentration tool is Lemma~\ref{lem:bounded-drift-concentration},
the appendix form of the Azuma--Hoeffding estimate~\cite{Azuma1967}.
For an adapted process $X_t$ with $|X_{t+1}-X_t|\le L$ for
$0\le t<T$, put
$d_t=\E[X_{t+1}-X_t\mid\mathcal F_t]$.  Its fixed-time conclusion is
\begin{equation}\label{sel:azuma}
 \Prb\left(\left|X_t-X_0-\sum_{j<t}d_j\right|\ge u\right)
 \le 2\exp\left(-\frac{u^2}{8TL^2}\right)
 \qquad(t\le T).
\end{equation}
The stopped auxiliary processes satisfy the same increment
hypothesis.  For domain sizes we will bound the absolute accumulated
drift; for degrees an upper bound suffices, and for occupancy we will
use a lower bound.

\paragraph{Domain sizes.}
For each initial slot $i$, begin an auxiliary process at $L_i(0)$
and track $L_i(t)/s(t)$ while the slot remains present.  In a
performed step, let $L_i^+$ be its size after conflicts and thinning
but before removing $G$.  This definition includes the hypothetical
update when $G=i$.  Equal survival in \eqref{sel:equal-survival} gives
\begin{equation}\label{sel:size-expectation}
 \E[L_i^+\mid\mathcal F_t]=(1-\alpha)L_i(t).
\end{equation}
If this domain's thinning count exceeds $q^\beta$, replace the
numerator $L_i^+$ by $L_i(t)$, update the auxiliary value to
$L_i(t)/s(t+1)$, and then stop it.  Otherwise update to the true
value $L_i^+/s(t+1)$ and stop after this update if $G=i$ or if the
whole process aborts.  Until such a stop, the auxiliary process
agrees with $L_i(t)/s(t)$.

The replacement changes the conditional expectation of the numerator
by at most $q\varepsilon_q$.  Consequently the conditional drift
of this rescaled track has absolute value at most
$q\varepsilon_q/s_{\min}$.  Its increments are $O(q^\beta)$:
there are at most $q^\beta$ conflict deletions in a slot by
\eqref{sel:graph-bounds}, and at most $q^\beta$ thinning deletions
on a step without replacement.  The scaling change contributes
only $O(1)$, by \eqref{sel:scale-lower}.  On a replaced step the
scaling change is the only contribution.

\paragraph{Remaining degrees.}
For each initial vertex $v$, begin at $D_v(0)$ and track
$D_v(t)/b(t)$ while $v$ is alive and its slot remains present.
The faster scale $b(t)$ accounts for the removal of the chosen slot.
To verify the required drift, let $\widetilde D_v$ be the true
updated neighbor count after a performed step, including removal of
$G$, whether or not $v$ itself survives.
For each currently counted neighbor $u$, its probability of
surviving before slot removal is $1-\alpha$.  If $u$'s own slot
is chosen, it cannot conflict with $W$, because there are no
within-slot edges.  Therefore
\[
 \Prb(u\hbox{ survives before removal and its slot is }G
       \mid\mathcal F_t)=(1-\lambda_u)/h(t)
       \ge(1-\alpha)/h(t).
\]
Subtracting this slot-removal probability from $1-\alpha$ and
summing over the currently counted neighbors gives
\begin{equation}\label{sel:degree-drift}
 \E[\widetilde D_v\mid\mathcal F_t]
 \le(1-\alpha)(1-1/h(t))D_v(t)
 =\frac{b(t+1)}{b(t)}D_v(t).
\end{equation}

We next bound the increments.  When $W\notin J(v)$, conflicts
remove at most $q^\beta$ counted neighbors by
\eqref{sel:exceptional}, and slot removal costs at most $q^\beta$
by \eqref{sel:graph-bounds}.  If instead $W\in J(v)$, or if the
thinning count in the current neighbor set exceeds $q^\beta$,
censor the auxiliary update: replace $\widetilde D_v$ by $D_v(t)$,
update to $D_v(t)/b(t+1)$, and then stop the track.  In every other
case update to $\widetilde D_v/b(t+1)$, and then stop if $v$ has
died, its slot has been removed, or the whole process has aborted.

The exceptional choice has conditional probability at most
\begin{equation}\label{sel:exception-probability}
 \Prb(W\in J(v)\mid\mathcal F_t)
 \le\frac{|J(v)|}{h(t)\min_iL_i(t)}
 \le\frac{2q^\beta}{a s_{\min}c_*q^2}
 =O(q^{\beta-2}).
\end{equation}
Since $D_v(t)\le C_*q$, censoring increases the numerator's
conditional expectation by at most
$O(q^{\beta-1})+C_*q\varepsilon_q$.  By
\eqref{sel:degree-drift} and \eqref{sel:scale-lower}, the rescaled
track thus has conditional drift at most $O(q^{\beta-1})$.
Its increments have absolute value $O(q^\beta)$.  Without
replacement the raw decrement is at most $3q^\beta$; with
replacement it is zero.  The relative change of $b(t)$ is
$O(1/q)$, uniformly because $h(t)\ge aq$, so its scaling
contribution is again $O(1)$.

This is the point at which we use the location of the exceptional
set $J(v)$, not just its size.  If $W\in J(v)$, then $W$ is adjacent
to $v$ or belongs to its slot.  After this step $v$ is therefore dead
or its slot has been removed.  A track stopped for this reason can
never be needed at a later barrier check.

\paragraph{Simultaneous estimates and absence of aborts.}
The domain and degree processes now have the bounded increments and
drift estimates needed for \eqref{sel:azuma}.  Apply that fixed-time
bound with deviation $q^{4/5}/2$ separately to each track and each
time.  There are $O(q^2)$ tracks and $O(q)$ times.
The individual exceptional probability is at most
\[
 2\exp\{-\Omega(q^{8/5-1-2\beta})\}
 =2\exp\{-\Omega(q^{7/20})\}.
\]
The total absolute drift of a domain-size track is $o(1)$ and
the total upper drift of a degree track is $O(q^\beta)$.
It follows, by a union bound and by $s(t),b(t)\le1$, that with
probability $1-o(1)$,
\begin{equation}\label{sel:concentration}
 |L_i(t)-s(t)L_i(0)|\le q^{4/5},\qquad
 D_v(t)\le b(t)D_v(0)+q^{4/5}
\end{equation}
at every time when the auxiliary values still agree with the true
rescaled quantities.

We can now rule out aborts.  Exclude the post-step abort event, whose
probability was already bounded by $o(1)$.  At a putative first
pre-step abort, every remaining slot and every alive vertex still has
a track agreeing with its true rescaled quantity.  A previous thinning
censorship would have caused a post-step abort, and a previous
exceptional choice would have killed the vertex or removed its slot.
Thus \eqref{sel:concentration} applies at this abort check itself.
Its first inequality rules out \eqref{sel:size-barrier} for large $q$.
The second gives
\begin{equation}\label{sel:hazard-barrier}
 r_v(t)
 \le \frac{s(t)(h(t)/n)C_*q+q^{4/5}}
              {h(t)(s(t)c_*q-q^{4/5})}
 =\frac{C_*}{c_*n}+o(q^{-1})<P_*/q.
\end{equation}
The error is uniform for $s(t)\ge s_{\min}$ and $h(t)\ge aq$.
Thus no pre-step abort occurs on this event.

To prepare the completion, make one final check at $t=t_{\max}$.
Abort unless the remaining minimum domain size $L_{\min}$ and maximum
degree $\Delta$ satisfy
\begin{equation}\label{sel:tail-check}
 L_{\min}\ge s(t_{\max})c_*q/2,
 \qquad \Delta/L_{\min}\le1/8.
\end{equation}
The same estimates imply
\[
 \frac{\Delta}{L_{\min}}
 \le\frac{h_0}{n}\frac{C_*}{c_*}+o(1)
 \longrightarrow \frac{aC_*}{2c_*}<\frac18.
\]
Therefore the main process reaches and passes its final check with
probability $1-o(1)$, uniformly in its input graph.

\subsection{Keeping points on every line}

The preceding estimates control the available vertices.  We next
control the points left on affine lines, under the additional
geometric hypotheses of Proposition~\ref{sel:selection}.
Fix a line $\ell$, start with $T_\ell$, and remove its intersections
with the main selected lines other than $\ell$.  Write $T(t)$ for
the number of points still present.  Each step decreases $T(t)$ by
at most one.

Consider a slot whose direction differs from that of $\ell$.
Conditional on the current history, at most $T(t)$ candidate lines
in that slot meet a remaining point: each such line is determined by
its intersection with $\ell$, and the assignment of lines within a
slot is injective.  A line in the same direction as $\ell$ makes no
deletion under our convention.  The size barrier therefore gives
\begin{equation}\label{sel:occupancy-drift}
 \Prb(T(t+1)=T(t)-1\mid\mathcal F_t)
 \le \frac{B_*T(t)}q,
 \qquad
 \E[T(t+1)\mid\mathcal F_t]\ge(1-B_*/q)T(t).
\end{equation}
Track $T(t)/(1-B_*/q)^t$, freezing on an abort and including a
performed step's increment before a post-step freeze.  For $t\le3q$
its denominator is at least $e^{-6B_*}$.  Its conditional drift
is nonnegative, and its increments have bounded absolute value
independently of $q$: the raw size changes by at most one and
$T(t)\le q$ controls the change of scale.

Apply the lower-tail consequence of \eqref{sel:azuma} to these
stopped processes and take a union bound over the $q(q+1)$ affine
lines.  The probability that the main stage completes while some
line has final main size below
\begin{equation}\label{sel:occupancy-main}
 c_0 e^{-6B_*}q/2=2\rho q
\end{equation}
is $o(1)$.
Any completion on the remaining $h_0$ slots removes at most $h_0$
further points from any one line.  By $a<\rho/2$ and
$h_0=\lceil aq\rceil$, the terminal sizes are at least $\rho q$
for large $q$.  Thus the required occupancy estimate will hold for
every successful tail completion on these histories.

\subsection{Completing the last slots}

After \eqref{sel:tail-check}, the remaining graph is sparse enough
for the Lov\'asz local lemma.  We need both an independent transversal
and a bound on prescribed choices under the law used to choose it.
The following conditional form supplies both, using the local-lemma
argument of Erd\H{o}s and Lov\'asz~\cite[Section~2]{ErdosLovasz1975}.

\begin{lemma}[Tail conditioning]\label{sel:tail-cylinder}
Let a graph be partitioned into slots of sizes $L_i$, with no
within-slot edges.  Let $L_{\min}=\min_iL_i$ and let its maximum
degree be $\Delta$.  If $L_{\min}\ge3$ and
$\Delta/L_{\min}\le1/8$, independent uniform choices from the
slots have positive probability of being an independent transversal.
Under the law conditioned on that event, a cylinder prescribing
values in a set $J$ of $s$ slots has probability at most
\begin{equation}\label{sel:tail-cylinder-bound}
 e^{s/2}\prod_{i\in J}L_i^{-1}.
\end{equation}
An impossible prescription has probability zero.
\end{lemma}

\begin{proof}
For every edge $e=uv$ between slots $i,j$, let $B_e$ be the event
that those two endpoints are chosen and put
$x_e=2/(L_iL_j)$.  Two events are neighbors when their edges use
a common slot.  Each $B_e$ is jointly independent of all events
outside its neighborhood.  For each slot $i$,
\[
 \sum_{e\text{ using }i}x_e
 \le\sum_{u\in Q_i}\frac{2\deg(u)}{L_iL_{\min}}
 \le\frac{2\Delta}{L_{\min}}\le\frac14.
\]
Thus a dependency neighborhood has total $x$-weight at most
$1/2$.  Since $x_e<1/2$,
\[
 \Prb(B_e)=x_e/2
 \le x_e\prod_{f\sim e}(1-x_f),
\]
using $\prod(1-x_f)\ge1-\sum x_f\ge1/2$.

We record the conditioning argument, since it also controls cylinders.
Induction on the size of an avoidance set $S$ of other events gives
positive avoidance probability and
\begin{equation}\label{sel:lll-induction}
 \Prb(B_e\mid\bigcap_{f\in S}B_f^c)\le x_e.
\end{equation}
Partition $S$ into nonneighbors $T$ and neighbors $U$ of $e$.
If $U$ is empty, joint independence gives the assertion.  Otherwise
the induction hypothesis, used successively on the events of $U$,
gives
\[
 \Prb\left(\bigcap_{f\in U}B_f^c
                  \,\middle|\,\bigcap_{f\in T}B_f^c\right)
 \ge\prod_{f\in U}(1-x_f).
\]
Every conditioning set in this product has smaller size than $S$.
The numerator after conditioning only on avoidance of $T$ is at most
$\Prb(B_e)$, by joint independence.  Dividing by the displayed
lower bound, and using
$\Prb(B_e)\le x_e\prod_{f\sim e}(1-x_f)$, proves
\eqref{sel:lll-induction}.  Positivity at each size follows by
successively avoiding its events and applying the preceding-size
bound.  Thus the induction establishes the probability bound and
positivity together.

Now fix a possible cylinder $C$ on slots $J$.  Divide the bad events
into those disjoint from $J$ and those using $J$.  Avoiding the former
leaves $\Prb(C)=\prod_{i\in J}L_i^{-1}$ unchanged, since those
events involve only other slots.  By \eqref{sel:lll-induction},
avoiding the latter increases the upper bound by at most their
product of $(1-x_e)^{-1}$.  Consequently
\[
 \Prb(C\mid\hbox{no bad event})
 \le\Prb(C)\prod_{e\text{ using }J}(1-x_e)^{-1}
 \le\Prb(C)\exp\left(2\sum_{e\text{ using }J}x_e\right)
 \le e^{s/2}\Prb(C).
\]
Here $-\log(1-x)\le2x$ for $0\le x\le1/2$ and the incident
weight per slot is at most $1/4$.  Counting an event twice in the
sum over prescribed slots only increases the upper bound.
\end{proof}

For every history passing \eqref{sel:tail-check}, choose the tail
according to the conditioned law in Lemma~\ref{sel:tail-cylinder}.
Its hypotheses hold for large $q$, since the check also gives
$L_{\min}\ge s_{\min}c_*q/2\ge3$.  The tail choices are mutually
compatible, and the earlier neighbor deletions make them compatible
with every main choice.  There is no further failure.  We have proved
\eqref{sel:success}; the occupancy estimate, which allowed any
successful completion of the last $h_0$ slots, also proves
\eqref{sel:occupancy-conclusion}.

\subsection{Cylinder probabilities for the whole selection}

The remaining claim is \eqref{sel:cylinder}.  We must bound the
probability of prescribed choices jointly with success, without
conditioning a main choice on the later success of the procedure.
For this purpose, realize the random slot order by sampling a uniform
permutation $\pi$ of all slots at the outset and using its first
$n-h_0$ entries.  Choose $\pi$ independently of all randomness for
vertex choices and thinning.  This is the same process as before:
the rates and abort decisions use no unexposed part of $\pi$, and
its next entry is uniform among the remaining slots in the original
filtration.

Only now, for the cylinder calculation, condition on the entire
permutation.  This divides the prescribed slots deterministically into
$s_{\rm m}$ main slots and $s_{\rm t}$ tail slots, with
$s_{\rm m}+s_{\rm t}=s$.  At any prescribed main slot reached before
abort, the vertex is still uniform in its current domain when $\pi$
is fixed.  The pre-step size barrier therefore bounds the probability
of its prescribed value by
\[
 \frac{2}{s_{\min}c_*q}.
\]
Multiplication of these conditional upper bounds at the specified
main steps bounds the probability of reaching and satisfying them
all.  An intervening abort only lowers that probability.  Given
any completed history passing the tail check,
Lemma~\ref{sel:tail-cylinder} bounds the tail prescription by
\[
 \left(\frac{2e^{1/2}}{s_{\min}c_*q}\right)^{s_{\rm t}}.
\]
Taking conditional expectations backwards from the tail and then
through the prescribed main steps gives
\[
 \Prb(\mathcal A,\hbox{ all prescriptions}\mid\pi)
 \le\left(\frac{2}{s_{\min}c_*q}\right)^{s_{\rm m}}
      \left(\frac{2e^{1/2}}{s_{\min}c_*q}\right)^{s_{\rm t}}.
\]
Average over $\pi$ and take
\begin{equation}\label{sel:cylinder-constant}
 K_* =\max\{1,2e^{1/2}/(s_{\min}c_*)\}.
\end{equation}
This proves \eqref{sel:cylinder} and completes
Proposition~\ref{sel:selection}.  The two uses of the slot order are
distinct: all drift estimates were taken in the original filtration,
where the next slot remained uniform.  Fixing the full permutation
was used only to multiply the cylinder bounds.  No drift identity
was asserted or used under that conditioning.

\subsection{Application to the deletion lines}\label{sec:selection-application}

We now apply the graph result to the geometric construction.
Proposition~\ref{prop:candidate-graph} supplies the domain and degree
bounds, the codegree bound, and the exceptional sets required by
Proposition~\ref{sel:selection}.  Its retained candidates are
injective in their lines within each slot, and an independent
transversal is a compatible collection of split tuples.  For the
occupancy statement, take
\[
 T_\ell=
 \begin{cases}
  \ell,&\ell\in\{M_g^1:g\in\mathcal G\},\\
  \ell\setminus\bigcup_g M_g^1,&\text{otherwise}.
 \end{cases}
\]
The same graph proposition gives $|T_\ell|\ge c_0q$.  Recall from
\eqref{eq:det-core} that $R_0$ avoids the deletion lines, whereas
$R$ avoids only the selected split lines.  If $\ell$ is neither a
selected split line nor one of the $M_g^1$, its final tracked set is
exactly $\ell\cap R_0$.  For $\ell=M_g^1$, the final set is
$\ell\cap R$, because a constituent cannot be a candidate split
line.  Hence \eqref{sel:occupancy-conclusion} gives the line
occupancy required by the deterministic construction.

To describe the resulting deletion-line law, include the constituent
sampling in the experiment.  Sample the two offsets in each direction
independently and with replacement, as in
Section~\ref{sec:candidates}.  Abort if the data fail the candidate
graph conditions or if any constituent pair coincides; otherwise run
the selection procedure.  Write $\mathcal A$ for the event that this
entire experiment returns a collection.  By
Proposition~\ref{prop:candidate-graph} and \eqref{sel:success}, there
is a fixed $a_0>0$ such that
\begin{equation}\label{sel:unconditional-success}
 \Prb(\mathcal A)\ge a_0
\end{equation}
for all sufficiently large $q$.  For every fixed offset outcome $M$,
\begin{equation}\label{sel:conditional-offset-cylinder}
 \Prb(\mathcal A,\hbox{ prescribed choices in }s\hbox{ slots}\mid M)
 \le(K_*/q)^s.
\end{equation}
These prescriptions may be chosen as functions of $M$, but are fixed
once $M$ is fixed.  If the initial check fails for that offset outcome,
the left side is zero.  This estimate lets us compare the law of the
returned deletion set with a uniform law without changing the initial
offset distribution.

\begin{corollary}[Domination of the deletion-line law]
\label{sel:deletion-domination}
Let $D$ consist of the selected split lines and the lines $M_g^1$
in the preceding experiment.  Conditional on $\mathcal A$, the law
of $D$ is dominated pointwise by $\exp(O(q))$ times the uniform
law on sets containing three distinct lines in each direction.
\end{corollary}

\begin{proof}
Fix such an unordered deletion set $D_0$.  For each direction there
are six assignments of its three lines to $M_g^1$, split slot 1,
and split slot 2.  For a fixed assignment in all directions, the
initial probability of the prescribed $M_g^1$ values is
$q^{-(q+1)}$.  Conditional on the full offset data, the probability
of all prescribed split lines jointly with $\mathcal A$ is at most
$(K_*/q)^{2(q+1)}$, by
\eqref{sel:conditional-offset-cylinder}.  Averaging over the
unconstrained offsets $M_g^0$ and summing over all role assignments
gives
\[
 \Prb(\mathcal A,D=D_0)
 \le\left(\frac{6K_*^2}{q^3}\right)^{q+1}.
\]
The uniform law $U$ assigns $D_0$ probability
$\binom q3^{-(q+1)}$.  Dividing also by
\eqref{sel:unconditional-success}, we obtain
\begin{equation}\label{sel:deletion-ratio}
 \frac{\Prb(D=D_0\mid\mathcal A)}{U(D_0)}
 \le a_0^{-1}\left[K_*^2(1-1/q)(1-2/q)\right]^{q+1}
 \le a_0^{-1}K_*^{2(q+1)}=\exp(O(q)).
\end{equation}
All probabilities before this final conditioning use the original
independent offset law; the event $\mathcal A$ enforces the candidate
graph conditions and constituent distinctness.
\end{proof}
